% ============================================================
%  2.5  Why the Map-Space Explodes
% ============================================================
\section{Why the Map-Space Explodes}
\label{sec:bg-mapspace}

Section~\ref{sec:bg-mapping} introduced the mapping problem and its
three degrees of freedom.
This section demonstrates that even modest workloads on relatively
simple architectures produce map-spaces of astronomical size, motivating
the need for analytical cost models.

Two observations underpin the explosion:

\begin{itemize}
  \item \textbf{Permutations per memory level.}
    Loop \emph{ordering} at memory levels matters because it determines
    which operand is stationary, whether halo reuse is available, and in
    what sequence transfers occur.
    In contrast, on spatial levels loop order does not change behavior:
    only the chosen \emph{parallelized dimensions} matter
    (Section~\ref{sec:loop-ordering}).

  \item \textbf{Independent tiling choices.}
    For each \emph{normal} dimension~$d \in D$
    (cf.\ Definition~\ref{def:normal-dim}), we decompose its extent
    into prime factors and distribute the resulting copies across the
    $L$ levels of the architecture.
    These choices are independent across distinct primes and across
    dimensions.
\end{itemize}


% ------------------------------------------------------------
\subsection{Formal Setup and Notation}
\label{sec:mapspace-notation}
% ------------------------------------------------------------

Let $L$ denote the total number of levels in the spatial architecture
and $L_m \le L$ the number of \emph{memory} levels.
Let $D$ be the set of normal kernel dimensions (affine dimensions
affect tile \emph{sizes} but are not independently looped; see
Section~\ref{sec:normal-affine}).
For each dimension $d \in D$, write the prime factorization of its
extent as
\[
  |d| \;=\; \prod_{p}\, p^{\,v_p(d)},
\]
where $v_p(d)$ is the $p$-adic valuation of~$|d|$ (i.e., the
exponent of prime~$p$ in the factorization).

\paragraph{Counting objective.}
We aim to count the total number of mappings \emph{ignoring validity
constraints}, that is, including configurations whose tiles may
exceed memory capacities.
This upper bound is useful to understand scale of the problem.


% ------------------------------------------------------------
\subsection{Counting Loop Permutations}
\label{sec:count-permutations}
% ------------------------------------------------------------

On each memory level, we may arbitrarily permute the $|D|$ loops.
This yields $|D|!$ choices per memory level and
\[
  (|D|!)^{L_m}
\]
in total.
Spatial levels do not contribute because, as established in
Section~\ref{sec:loop-ordering}, loop order on them is irrelevant, so no permutation is added.


% ------------------------------------------------------------
\subsection{Counting Factor Allocations}
\label{sec:count-factors}
% ------------------------------------------------------------

For a fixed dimension~$d$ and a fixed prime~$p$, we must distribute
$v_p(d)$ indistinguishable copies of~$p$ across $L$ \emph{labeled}
levels; levels may receive zero copies, and the assignment to
different levels produces different tile sizes, so order matters.
The number of such allocations is the number of weak compositions
of~$v_p(d)$ into~$L$ parts:
\[
  \binom{v_p(d) + L - 1}{L - 1}.
\]

Because allocations are independent across distinct primes $p \mid d$
and across dimensions $d \in D$, the \emph{total} number of factor
allocations is
\[
  \prod_{d \in D}\;\prod_{p \mid d}
    \binom{v_p(d) + L - 1}{L - 1}.
\]


% ------------------------------------------------------------
\subsection{Map-Space Cardinality}
\label{sec:mapspace-card}
% ------------------------------------------------------------

Combining permutations and factor allocations yields the following
expression for the (valid + invalid) map-space
size~\cite{ronzani2025quickflow}:

\begin{equation}
\label{eq:ms-cardinality}
\boxed{%
  |\mathcal{MS}|
  \;=\;
  \underbrace{(|D|!)^{L_m}}_{\text{loop permutations}}
  \;\cdot\;
  \underbrace{
    \prod_{d \in D}\;\prod_{p \mid d}
      \binom{v_p(d) + L - 1}{L - 1}
  }_{\text{factor allocations}}.
}
\end{equation}

\paragraph{Growth with architecture complexity.}
The permutation term scales \emph{exponentially} in~$L_m$ (with
base~$|D|!$), while the factor-allocation term grows
\emph{polynomially} in~$L$, with the exponent equal to the total
number of prime copies
$\sum_{d \in D}\sum_{p \mid d} v_p(d)$,
which itself increases with~$|D|$ and the loop extents.


% ------------------------------------------------------------
\subsection{Worked Example}
\label{sec:mapspace-example}
% ------------------------------------------------------------

Consider a standard 2D convolution with $|D| = 6$ normal loops
(e.g., $M, C, P, Q, R, S$) mapped onto an architecture with $L = 5$
levels, of which $L_m = 4$ are memory levels.

\paragraph{Permutations.}
\[
  \text{Permutations}
  \;=\; (6!)^{4}
  \;=\; 720^{4}
  \;\approx\; 2.68 \times 10^{11}.
\]

\paragraph{Factor allocations.}
Suppose each dimension has, in total, five prime-factor copies to
distribute across the five levels.

Two boundary cases illustrate the range:

\begin{enumerate}
  \item \textbf{All five copies on the same prime}
    (e.g., $v_p(d) = 5$ unique prime):
for each dimension
    $\binom{5 + 5 - 1}{5 - 1} = \binom{9}{4} = 126$,
    so across six dimensions $126^{6} \approx 1.07 \times 10^{13}$.

  \item \textbf{Five distinct primes, one copy each}
    ($v_{p}(d) = 1$ for five primes):
    per dimension
    $\bigl(\binom{1 + 5 - 1}{5 - 1}\bigr)^{5} = 5^{5} = 3125$,
    so across six dimensions $3125^{6} \approx 3.05 \times 10^{20}$.
\end{enumerate}

\paragraph{Total.}
Multiplying by permutations, the map-space size ranges between
\[
  \underbrace{720^{4} \cdot 126^{6}}_{\approx\;1.1 \times 10^{24}}
  \quad\text{and}\quad
  \underbrace{720^{4} \cdot 3125^{6}}_{\approx\;2.5 \times 10^{32}}.
\]

Even the conservative end of this range \emph{far} exceeds
$10^{24}$.
These numbers represent \emph{all} mappings, including invalid ones;
nevertheless, they demonstrate that exhaustive enumeration is
infeasible.

\paragraph{Practical implications.}
Across realistic map-spaces, EDP can vary by up to $10^{4}\times$
between the best and worst valid mappings, and near-optimal points are
sparsely distributed~\cite{kao2020gamma}.
Small random sampling budgets are therefore unlikely to discover good
solutions without \emph{structural pruning}, heuristics that
eliminate provably sub-optimal or redundant regions of the search
space before evaluation.
This motivates the use of analytical cost models that can evaluate
candidate mappings cheaply (Section~\ref{sec:bg-cost-model}) and
mapper algorithms that exploit the structure of the map-space to
navigate it efficiently.